\section{Conception et Architecture Réseau}

%------------------------------------------------------------
% Slide 7 : Objectifs de la Conception
%------------------------------------------------------------
\begin{frame}{Objectifs de la Conception Réseau}
    Pour sécuriser l'institut \textit{DataLab Research}, la conception repose sur deux piliers :
    
    \vspace{0.3cm}
    \begin{block}{Segmentation par VLAN (IEEE 802.1Q)}
        \begin{itemize}
            \item Réduction drastique des domaines de diffusion.
            \item Isolation logique complète des différentes catégories d'utilisateurs.
            \item Application de politiques de filtrage strictes à la frontière des zones.
        \end{itemize}
    \end{block}

    \begin{block}{Cœur de Réseau Centralisé}
        Mise en œuvre d'un routage inter-VLAN de type \textbf{Router-on-a-Stick} pour interconnecter de manière contrôlée tous les réseaux de l'institut.
    \end{block}
\end{frame}

%------------------------------------------------------------
% Slide 8 : Plan de Segmentation Exhaustif
%------------------------------------------------------------
\begin{frame}{Plan de Segmentation Exhaustif (VLANs)}
    L'ingénierie globale cible une sectorisation stricte en \textbf{6 zones distinctes} :
    
    \vspace{0.2cm}
    \begin{columns}[T]
        \begin{column}{0.48\textwidth}
            \begin{itemize}
                \item \textbf{VLAN 10 (DMZ) :} Zone publique tampon exposée à l'extérieur.
                \item \textbf{VLAN 20 (Production) :} Hébergement du cœur des services applicatifs.
                \item \textbf{VLAN 30 (Chercheurs) :} Dédié aux stations des 80 scientifiques.
            \end{itemize}
        \end{column}
        
        \begin{column}{0.48\textwidth}
            \begin{itemize}
                \item \textbf{VLAN 40 (Ingénieurs) :} Segment isolé pour le développement.
                \item \textbf{VLAN 50 (Supervision) :} Flux d'administration et métriques système.
                \item \textbf{VLAN 60 (Backup) :} Zone ultra-sécurisée dédiée au stockage répliqué.
            \end{itemize}
        \end{column}
    \end{columns}
\end{frame}

%------------------------------------------------------------
% Slide 9 : Plan d'Adressage IP Général
%------------------------------------------------------------
\begin{frame}{Plan d'Adressage IP Implémenté}
    Utilisation du bloc d'adresses privé global \texttt{10.0.0.0/16}, découpé en sous-réseaux de classe C (\texttt{/24}).

    \vspace{0.2cm}
    \centering
    \small
    \begin{tabular}{llll}
        \toprule
        \textbf{VLAN / Zone} & \textbf{Plage Réseau} & \textbf{Masque} & \textbf{Passerelle (Gateway)} \\
        \midrule
        \textbf{VLAN 10} -- DMZ & \texttt{10.0.10.0/24} & \texttt{255.255.255.0} & \texttt{10.0.10.254} \\
        \textbf{VLAN 20} -- Production & \texttt{10.0.20.0/24} & \texttt{255.255.255.0} & \texttt{10.0.20.254} \\
        \textbf{VLAN 30} -- Chercheurs & \texttt{10.0.30.0/24} & \texttt{255.255.255.0} & \texttt{10.0.30.254} \\
        \textbf{VLAN 60} -- Backup & \texttt{10.0.60.0/24} & \texttt{255.255.255.0} & \texttt{10.0.60.254} \\
        \bottomrule
    \end{tabular}
    
    \vspace{0.3cm}
    \flushleft
    \footnotesize{\textbf{Note d'ingénierie :} Ce découpage garantit jusqu'à 254 hôtes par segment et offre une excellente lisibilité lors des analyses de paquets.}
\end{frame}

%------------------------------------------------------------
% Slide 10 : Topologie Générale de la Maquette GNS3
%------------------------------------------------------------
\begin{frame}{Topologie Générale sous GNS3}
    \begin{columns}[T]
        \begin{column}{0.4\textwidth}
            \textbf{Description de la Maquette :}
            \begin{itemize}
                \item Émulation réseau complète via le simulateur \textbf{GNS3}.
                \item Commutateur gérant le marquage et le transit des trames \textit{Trunk}.
                \item Connexion vers le routeur cœur de réseau Cisco.
                \item Intégration du serveur de production principal et du nœud de sauvegarde.
            \end{itemize}
        \end{column}
        
        \begin{column}{0.6\textwidth}
            \begin{figure}
                \centering
                % Intégration de ta capture d'architecture globale
                \includegraphics[width=0.98\textwidth]{figures/archi_globale.jpeg}
                \caption{Topologie logique et interconnexions émulées}
            \end{figure}
        \end{column}
    \end{columns}
\end{frame}